\section*{Acknowledgements}
This paper makes use of the following ALMA data: the \NProjCodes{} projects
listed in Table~\ref{tab:projcodes}, each cited as \texttt{ADS/JAO.ALMA\#}
followed by its project code. The \NEB{} primary-census execution blocks
searched span \PrEbDateFirst{} to \PrEbDateLast{}, and we thank the principal
investigators and teams who proposed and obtained
them.\ProjCodeUnresClause{} ALMA is a partnership of ESO
(representing its member states), NSF (USA) and NINS (Japan), together with NRC
(Canada), NSTC and ASIAA (Taiwan), and KASI (Republic of Korea), in cooperation
with the Republic of Chile. The Joint ALMA Observatory is operated by ESO,
AUI/NRAO and NAOJ. This work has made use of data from the European Space
Agency (ESA) mission {\it Gaia} \citep{GaiaDR3}, processed by the {\it Gaia}
Data Processing and Analysis Consortium (DPAC). This research has made use
of the SIMBAD database \citep{SIMBAD2000} (CDS, Strasbourg, France); of the
NASA Exoplanet Archive \citep{NEA2013}, which is operated by the California
Institute of Technology under contract with the National Aeronautics and
Space Administration under the Exoplanet Exploration Program; and of the
Splatalogue spectral-line database \citep{Splatalogue2007}.

\emph{Use of generative artificial intelligence.} Large language models were
used in this work. They wrote most of the software: the archival retrieval and
calibration scripts, the search and vetting pipeline, the statistical tests,
and the generators that produce every number, table and figure in this paper.
They also drafted text, which the authors revised. The authors set the
scientific question, chose every criterion, and are responsible for the
content of this paper.

\emph{Testing and validation of that analysis.} No result here is accepted
on a model's assertion. No number in this paper is typed: each is written by
the program that computes it from the data, so a value in the text cannot
differ from the analysis it cites -- which is not the same as saying that no
other extraction of the same data would give a different one, and
Table~\ref{tab:confirm} names the one window where the archive holds two
extractions that differ. A separate audit re-derives the
identities that must hold among the published values -- counts that must
partition, quantities that must sum, numbers reached by two routes that must
agree -- and refuses the paper where one does not. Every assertion in the
analysis was first made to fail on purpose, on deliberately wrong input,
before it was relied on, since an assertion that has never fired is no
evidence that anything is right. Every reference cited here has been checked
against its source of record.

\emph{Software.} The pipeline depends on \textsc{casa}~6.7.6
\citep{CASA2022}, \textsc{python}~3 with \textsc{astropy} \citep{Astropy2022}, \textsc{numpy}
\citep{Harris2020} and \textsc{scipy} \citep{Virtanen2020}, and the ALMA Science
Archive and SIMBAD TAP interfaces through \textsc{astroquery} \citep{Ginsburg2019};
all are free and open-source, at the versions cited.

\emph{Author contributions.} G.J.W.\ conceived the survey, defined the target sample and the selection
criteria, and wrote the manuscript. R.D.\ built and operated the archival
retrieval, calibration and search pipeline, implemented the spatial control
statistic and the injection--recovery campaigns, and produced the per-window
catalogue and its regeneration scripts. Both authors designed the vetting
procedure, chose the masked species, carried out the statistical
validation reported in \S\ref{sec:statistic} and
Appendix~\ref{app:falsealarm}, and agreed the dispositions of
the crossing ledger.
%% Authors to confirm this division before submission.

\emph{Funding.} This work received no dedicated funding. It used public archival data and
no observing time was awarded for it.
%% Authors to replace this if any grant supported the work.

\emph{Competing interests.} The authors declare no competing interests. VBRL Holdings is a privately
held company through which R.D.\ conducts independent research, named
here solely as his affiliation; it had no role in the design of the
survey, the selection of targets, the analysis, the decision to publish
or the content of this paper, and provided no funding for the work.
%% Authors to confirm this wording, which is accurate only if the company's
%% involvement was indeed nil.

\section*{Data availability}
\label{sec:data}
\emph{Public.} Every observation this paper rests on is public and is
identified here precisely enough to be fetched. The \NEB{} primary-census
execution blocks, and those of every other set of
Table~\ref{tab:blockledger}, are held in the ALMA Science Archive under the
\NProjCodes{} project codes of Table~\ref{tab:projcodes} and are retrieved by
the execution-block identifier, which the crossing ledger gives for each of
the \EvNCrossRaw{} threshold crossings; archival data should carry the full
provenance of their origin and the credit for it \citep{Chen2022}.
The astrometry is Gaia~DR3, the planet counts the NASA Exoplanet Archive, and
the mask one Splatalogue query of CDMS and JPL made on \McQueryDate{} and
repeatable as it stands.

Five plain-text files accompany the paper on the arXiv, served beside the
article and not part of it, each describing its own columns in its own header:
the per-window catalogue, one row for each of the \NWindows{} searched windows,
with \SiCatColList; the crossing ledger, one row for each of the
\EvNCrossRaw{} crossings, with \SiLdgColList{} and sampled in
Table~\ref{tab:ledger}; the master counts, all
\MtCountRows{} rows of Table~\ref{tab:counts}; for each crossing carrying
a quoted exclusion, the probability that a carrier at the trigger amplitude
would have satisfied the recurrence criterion, calculated and measured; and,
for the same crossings, the exclusion depth read over separated blocks alone
and over the matched registration tolerance.

\emph{What they support.} With the tables here -- Table~\ref{tab:blockledger}
for the block and window partition, Table~\ref{tab:counts} for the headline
counts, and Appendix~\ref{app:persys} for every system of the carrier search
-- they rebuild the window, star
and system counts; the crossing counts and their classification; the chance
expectation of \CeExp{} crossings, as a sum over each Class~A window's own
control exceedance rate and masked share; the $\mathrm{EIRP}_{90}$ summaries,
including which windows carry an injected completeness and which a transferred
one; and every exclusion depth with its recovery probability. Two quantities
they do not rebuild: the sub-trigger stellar-to-control step of
Appendix~\ref{app:falsealarm}, measured on above-threshold cell counts that are
not released, and the interval on the chance expectation, whose resampling the
per-window rates support only in part.

\emph{Not public.} The calibrated visibilities, the per-window products and the
code that produced them are not distributed, so this execution of the search
cannot be re-run; the method is specified in enough detail to be
reimplemented, which is the stronger test in any case.
